\documentclass[11pt]{article}
\usepackage{amsmath,amssymb,amsthm}
\usepackage{geometry}
\geometry{margin=1in}
\newtheorem{lemma}{Lemma}
\newtheorem{theorem}{Theorem}
\theoremstyle{remark}
\newtheorem{remark}{Remark}

\title{Tangential Eigenvalue Universality on Constant-Curvature Surfaces}
\author{Sandbox for Paper~I, Item~\#2}
\date{\today}

\begin{document}
\maketitle

\section*{Main result}

\begin{theorem}[Riemannian Havelock universality]\label{thm:rh-univ}
On the unit $2$-sphere, the tangential eigenvalue of the constrained
Hessian $\nabla^2(H_{\mathrm{sph}} - \Omega J)$ at mode
$m \in \{1, \ldots, N-1\}$ is
\[
  \hat e_t^{\,T}\,\nabla^2(H_{\mathrm{sph}} - \Omega J)\,\hat e_t
  \;=\; \frac{m(N-m)}{2\xi},
\]
identical to the flat-plane result at ring radius $\sqrt{\xi}$.
The same holds on the hyperbolic plane with $(1+\xi)$ replaced by
$(1-\xi)$ throughout.
\end{theorem}

\begin{proof}
Decompose $H_{\mathrm{sph}} = H_{\mathrm{flat}} + W$ where
$H_{\mathrm{flat}} = -\sum_{j<k}\ln|z_j - z_k|$ and
$W = \tfrac{N-1}{2}\sum_k \ln(1+|z_k|^2)$.  The Lagrangian is
$\mathcal{L} = H_{\mathrm{flat}} + W - \Omega J$ with
$J = \sum_k(1-|z_k|^2)/(1+|z_k|^2)$.  The tangential eigenvalue
decomposes as
\[
  \hat e_t^T \nabla^2\mathcal{L}\,\hat e_t
  = \underbrace{\hat e_t^T \nabla^2 H_{\mathrm{flat}}\,\hat e_t}_{
      \text{flat Havelock}}
  + \underbrace{\hat e_t^T \nabla^2(W - \Omega J)\,\hat e_t}_{
      \text{curvature + constraint}}.
\]
We show the second term equals $(N-1)/(2\xi)$, i.e.\ the flat-plane
Lagrange shift $-2\mu_L = (N-1)/(2\xi)$ at ring radius $\sqrt{\xi}$.
Combined with the flat-plane unshifted tangential eigenvalue
$(m(N-m) - (N-1))/(2\xi)$ (derived in the constrained-minimum proof,
\S6.2), this gives $m(N-m)/(2\xi)$.

\medskip
\textbf{Step 1: Tangential projection of $\nabla^2 W$.}
At vertex $k$ on the ring ($|z_k|^2 = \xi$, angle $\theta_k$), the
Hessian of $W_k = \tfrac{N-1}{2}\ln(1+|z_k|^2)$ has entries
\begin{align*}
  \frac{\partial^2 W_k}{\partial x_k^2}
    &= \frac{N-1}{1+\xi} - \frac{2(N-1)\xi\cos^2\!\theta_k}{(1+\xi)^2},
  &
  \frac{\partial^2 W_k}{\partial y_k^2}
    &= \frac{N-1}{1+\xi} - \frac{2(N-1)\xi\sin^2\!\theta_k}{(1+\xi)^2},
  \\
  \frac{\partial^2 W_k}{\partial x_k\partial y_k}
    &= -\frac{2(N-1)\xi\sin\theta_k\cos\theta_k}{(1+\xi)^2}.
\end{align*}
The mode-$m$ tangential unit vector at vertex~$k$ is
$\hat e_{t,k} \propto \alpha_k(-\sin\theta_k, \cos\theta_k)$
with $\alpha_k = \cos(2\pi mk/N)$.  The tangential projection
at vertex~$k$ is
\begin{align*}
  \sin^2\!\theta_k\,\frac{\partial^2 W_k}{\partial x_k^2}
  &- 2\sin\theta_k\cos\theta_k\,\frac{\partial^2 W_k}{\partial x_k\partial y_k}
  + \cos^2\!\theta_k\,\frac{\partial^2 W_k}{\partial y_k^2} \\
  &= \frac{N-1}{1+\xi}
    - \frac{2(N-1)\xi}{(1+\xi)^2}
      \bigl[\sin^2\!\theta_k\cos^2\!\theta_k
            - 2\sin^2\!\theta_k\cos^2\!\theta_k
            + \cos^2\!\theta_k\sin^2\!\theta_k\bigr] \\
  &= \frac{N-1}{1+\xi}.
\end{align*}
The bracketed expression vanishes identically ($1 - 2 + 1 = 0$), so
the projection is $\theta_k$-independent and hence \emph{mode-independent}.
Summing over vertices with $\sum_k \alpha_k^2 = N/2 = \|\hat e_t\|^2$:
\begin{equation}\label{eq:W-tt}
  \hat e_t^T\,\nabla^2 W\,\hat e_t = \frac{N-1}{1+\xi}.
\end{equation}

\textbf{Step 2: Tangential projection of $\nabla^2 J$.}
With $J_k = (1 - |z_k|^2)/(1 + |z_k|^2)$:
\begin{align*}
  \frac{\partial^2 J_k}{\partial x_k^2}
    &= -\frac{4}{(1+\xi)^2} + \frac{16\xi\cos^2\!\theta_k}{(1+\xi)^3},
  &
  \frac{\partial^2 J_k}{\partial y_k^2}
    &= -\frac{4}{(1+\xi)^2} + \frac{16\xi\sin^2\!\theta_k}{(1+\xi)^3},
  \\
  \frac{\partial^2 J_k}{\partial x_k\partial y_k}
    &= \frac{16\xi\sin\theta_k\cos\theta_k}{(1+\xi)^3}.
\end{align*}
The same tangential projection yields
\[
  -\frac{4}{(1+\xi)^2}
  + \frac{16\xi}{(1+\xi)^3}\bigl[\sin^2\!\theta\cos^2\!\theta
    - 2\sin^2\!\theta\cos^2\!\theta + \cos^2\!\theta\sin^2\!\theta\bigr]
  = -\frac{4}{(1+\xi)^2},
\]
again $\theta$-independent.  Therefore
\begin{equation}\label{eq:J-tt}
  \hat e_t^T\,\nabla^2 J\,\hat e_t = -\frac{4}{(1+\xi)^2}.
\end{equation}

\textbf{Step 3: Combining $W - \Omega J$.}
With $\Omega = (N-1)(1-\xi^2)/(8\xi) = (N-1)(1-\xi)(1+\xi)/(8\xi)$:
\begin{align*}
  \hat e_t^T\,\nabla^2(W - \Omega J)\,\hat e_t
  &= \frac{N-1}{1+\xi}
     + \frac{(N-1)(1-\xi)(1+\xi)}{8\xi}\cdot\frac{4}{(1+\xi)^2} \\
  &= \frac{N-1}{1+\xi}
     + \frac{(N-1)(1-\xi)}{2\xi(1+\xi)} \\
  &= \frac{N-1}{1+\xi}\Bigl[1 + \frac{1-\xi}{2\xi}\Bigr]
   = \frac{N-1}{1+\xi}\cdot\frac{1+\xi}{2\xi}
   = \frac{N-1}{2\xi}.
\end{align*}
This equals exactly $-2\mu_L = (N-1)/(2\xi)$, the flat-plane Lagrange
shift at ring radius $\sqrt{\xi}$.

\textbf{Step 4: Assembly.}
The flat-plane unshifted tangential eigenvalue at radius $\sqrt{\xi}$ is
$(m(N-m) - (N-1))/(2\xi)$ (from the constrained-minimum derivation).
Adding the curvature-plus-constraint contribution $(N-1)/(2\xi)$:
\[
  \hat e_t^T\,\nabla^2\mathcal{L}\,\hat e_t
  = \frac{m(N-m) - (N-1)}{2\xi} + \frac{N-1}{2\xi}
  = \frac{m(N-m)}{2\xi}.
\]

\textbf{Hyperbolic case.}
On $\mathbf{H}^2$ with $W = +(N{-}1)/2\sum_k\ln(1-|z_k|^2)$ and
$J = \sum_k(1+|z_k|^2)/(1-|z_k|^2)$ (the Poincar\'e-disk
formulation), the same tangential-projection argument gives
$\hat e_t^T\nabla^2 W\,\hat e_t = -(N{-}1)/(1-\xi)$ and
$\hat e_t^T\nabla^2 J\,\hat e_t = +4/(1-\xi)^2$, and with
$\Omega_{\mathbf{H}^2} = -(N{-}1)(1+\xi)(1-\xi)/(8\xi)$ the
combination gives
\[
  -\frac{N{-}1}{1{-}\xi}
  + \frac{(N{-}1)(1{+}\xi)(1{-}\xi)}{8\xi}\cdot\frac{4}{(1{-}\xi)^2}
  = -\frac{N{-}1}{1{-}\xi} + \frac{(N{-}1)(1{+}\xi)}{2\xi(1{-}\xi)}
  = \frac{N{-}1}{2\xi}.
\]
\end{proof}

\section{Numerical verification}

All intermediate identities are verified in
\texttt{numerical\_check.py} (this directory):
\begin{itemize}
\item $\hat e_t^T\nabla^2 W\,\hat e_t = (N-1)/(1+\xi)$: confirmed
  numerically for $N \in \{5,6,7\}$, $\xi \in \{0.2, 0.5\}$, all
  modes $m \in \{2,\ldots,N-2\}$ (mode-independent to $10^{-2}$).
\item $\hat e_t^T\nabla^2 J\,\hat e_t = -4/(1+\xi)^2$: same
  parameters, mode-independent to $10^{-2}$.
\item $\hat e_t^T\nabla^2(W - \Omega J)\,\hat e_t = (N-1)/(2\xi)$:
  confirmed for $N \in \{5,6,7\}$, $\xi \in \{0.15, 0.3, 0.5, 0.7\}$.
\item Full tangential eigenvalue $= m(N-m)/(2\xi)$: confirmed for
  $N \in \{4,5,6,7\}$, $\xi \in \{0.15, 0.3, 0.5, 0.7\}$, all modes.
\item Symbolic: the cross-term cancellation
  $\sin^2\!\theta\cos^2\!\theta - 2\sin^2\!\theta\cos^2\!\theta
  + \cos^2\!\theta\sin^2\!\theta = 0$ is verified by \texttt{sympy}
  in the explicit tangential projection of $\nabla^2 W$.
\end{itemize}

\end{document}
